\section{Related Work}
\label{sec:related_work}

In this section, we situate our work within the broader context of model merging, test-time adaptation, and the recent critical audits of merging complexity.

\subsection{Static Model Merging}
Static model merging methods aim to combine multiple fine-tuned models in a single forward pass without further training or optimization. The foundational baseline, Task Arithmetic \cite{ilharco2022editing}, creates task vectors by subtracting pre-trained weights from fine-tuned weights and adds them linearly to the pre-trained base model. While elegant, task arithmetic suffers from destructive parameter interference when combining multiple, disparate tasks. To mitigate this, TIES-Merging \cite{yadav2023ties} resolves weight conflicts by pruning small-magnitude updates and resolving sign consensus. Similarly, DARE \cite{yu2023dare} employs random parameter dropping and scaling to sparsify updates, while ZipIt! \cite{stoica2023zip} merges models across different basins by aligning neuron correspondences. Other static methods weight parameters according to Fisher information \cite{matena2021merging} or compute closed-form least-squares linear mappings \cite{choshen2022fusing}. While these approaches are fully training-free and require zero parameters, their static nature limits their capability to adapt dynamically to specific test-time streams, often leading to sub-optimal performance under domain shift or complex task distributions.

\subsection{Test-Time Adaptive Model Merging}
To overcome the limitations of static merging, test-time adaptive model merging uses unlabeled target streams to dynamically optimize merging parameters. This field draws inspiration from classical test-time adaptation (TTA) techniques, such as TENT \cite{wang2020tent}, SHOT \cite{liang2020shot}, and MEMO \cite{zhang2022memo}. In the context of merging, AdaMerging \cite{yang2024adamerging} was among the first to introduce test-time adaptation of merging coefficients. AdaMerging parameterizes merging with both task-wise and layer-wise coefficients (amounting to over a thousand parameters for typical vision transformers) and optimizes them via entropy minimization or teacher-expert KL divergence. SyMerge \cite{jung2025symerge} extends this by optimizing low-rank coordinate-scaling adapters to dynamically realign feature spaces. FoldMerge \cite{foldmerge} escalates this approach by training a 4-layer RealNVP normalizing flow network with 2.6 million parameters to learn a non-linear weight-space diffeomorphism, warping the parameters of expert teachers into a shared alignment space.

While these adaptive methods achieve superior average accuracy, they introduce substantial parameter overhead (e.g., millions of adapter parameters) and require significant computational and memory footprint during test-time adaptation. Crucially, as we discuss below, their large parameter capacity exposes them to transductive overfitting.

\subsection{Critical Audits of Merging Complexity}
Recently, a series of critical auditing papers has begun to question the necessity of high-capacity and complex post-hoc model-merging pipelines, advocating for simpler, more robust alternatives:
\begin{itemize}
    \item {\bf Sharpness-Awareness Deconstruction:} The SAIM Audit \cite{saimaudit} methodologically deconstructs Sharpness-Aware Isotropic Merging (SAIM). The authors demonstrate that training expert models using sharpness-aware minimization (SAM) \cite{foret2020sharpness} is the fundamental driver of successful merging, rendering complex post-hoc SVD-based isotropic manipulation pipelines largely redundant. This supports our minimalist philosophy that proper training/flatness of the experts is far more critical than post-hoc overengineering.
    \item {\bf Overfitting of Layer-wise Parameters:} The Layer-wise Model Merging Sanity Check \cite{layerwiseaudit} deconstructs the assumption that fine-grained, layer-wise merging coefficients are necessary to resolve task weight interference. They prove that unconstrained optimization of layer-wise parameters (e.g., 52 parameters) on small calibration streams is highly prone to transductive overfitting under first-order gradient descent. Spatially averaging these optimized parameters (reducing parameter counts by 92.3\%) acts as a powerful spatial regularizer, yielding better generalization on unseen out-of-distribution streams.
\end{itemize}

Our proposed method, BPAM, directly builds upon these critical insights. Instead of scaling up parameter dimensions or training complex flow-based neural networks, we aggressively prune the trainable parameters to exactly $K$ global scalars. We leverage physical, geometric constraints (convexity and proximity regularization) to anchor the optimization, achieving a mathematically elegant and highly robust merging framework that honors Occam's razor.
